\documentclass[11pt]{article}

\usepackage[a4paper,margin=1in]{geometry}
\usepackage{amsmath,amssymb,amsthm}
\usepackage{booktabs}
\usepackage{multirow}
\usepackage{array}
\usepackage{graphicx}
\usepackage{hyperref}
\usepackage{xcolor}
\usepackage{microtype}
\usepackage{fancyvrb}
\usepackage{enumitem}
\usepackage[numbers,square]{natbib}
\usepackage{caption}
\usepackage{tabularx}
\usepackage{subcaption}

\hypersetup{
  colorlinks=true,
  linkcolor=blue!50!black,
  citecolor=blue!50!black,
  urlcolor=blue!50!black,
}

\newcommand{\hurts}{\text{$\dagger$}}
\newcommand{\helps}{\text{$\star$}}
\newcommand{\imhyp}{\textsc{im-Hyp}}
\newcommand{\imcot}{\textsc{im-CoT}}
\newcommand{\reactive}{\textsc{Reactive}}
\newcommand{\errk}{E_{\text{asym}}}

\title{Apply, Don't Articulate:\\
A Multi-Paradigm, Multi-Scale Study of\\
Internal-Model Prompting in LLMs}

\author{Mohamed A. Mabrok\\
\href{mailto:m.a.mabrok@gmail.com}{m.a.mabrok@gmail.com}}

\date{April 2026}

\begin{document}

\maketitle

\begin{abstract}
The Internal Model Principle (IMP) of \citet{francis1976internal}
states that any feedback controller that achieves perfect
asymptotic tracking of a class of exogenous signals must reproduce,
as an internal subsystem, the dynamics that generate those
signals.  A practitioner reading of this principle, popular in
agentic LLM design, prescribes that prompting a large language
model to \emph{articulate} an internal model of its environment
before acting should improve closed-loop performance, increasingly
so with scale.  We test this prescription in two distinct symbolic
paradigms: deterministic integer sequence prediction (a discovery
benchmark of $30$ hand-curated tasks plus a confirmation benchmark
of $260$ programmatically-generated tasks with multi-seed sampled
decoding) and code debugging ($30$ buggy Python functions with
hidden test-suite grading).  We apply both benchmarks across up to
five model scales (Qwen2.5-7B, 14B, 32B; Llama3.1-70B 4-bit;
Qwen2.5-Coder-7B, 14B, 32B), for a total of over $10{,}000$
prompted trials.  Across this corpus, four findings stand
together.  \emph{First}, articulating an internal model never
significantly helps in either paradigm, at any scale, under any
of three articulation phrasings tested.  \emph{Second}, the
IM-Hypothesis prompt significantly hurts at $22$ of $24$
difficulty $\times$ scale $\times$ phrasing cells in the
sequence-prediction confirmation benchmark, and produces a
\textbf{$97$-percentage-point pass-rate collapse} at Qwen2.5-Coder-14B
on the code-debugging benchmark ($t = 25$--$42$, $p < 10^{-20}$
at every difficulty band) -- the largest single prompt-induced
regression observed in the study.  \emph{Third}, the harm
exhibits a scale-dependent inverted-$U$ in both paradigms,
peaking at $14$--$32$B and weakening at the smallest and largest
tested scales.  \emph{Fourth}, four causal-mechanism arms
(self-consistency, oracle, hypothesis-first, equal-budget),
applied in both paradigms, refute the simplest tokenized-commitment
account but support a sharper one: \emph{the harm is in
rule-generation, not rule-application}.  Self-consistency over
$5$ sampled diagnoses does not recover the reactive baseline at
the catastrophe scale; reordering the prefix does not erase the
harm; equal-budget is inert.  But \emph{providing} the rule
directly (the oracle arm) helps significantly with growing
magnitude on hard tasks ($-0.84$ log-error gap at Qwen2.5-32B
on hard sequences; $+0.10$ pass-rate gain at Qwen2.5-Coder-32B on
hard code).  These results yield a clean deployment recommendation
that survives across two paradigms and seven model scales:
\emph{do not ask the LLM to articulate its internal model;
provide the model and ask it to apply.}  The benchmarks, prompts,
graders, and analysis pipeline are released.

\bigskip
\noindent\textbf{Keywords:} internal model principle, large
language models, sequence prediction, code debugging, prompt
engineering, chain-of-thought, control theory, scaling, negative
result, reproducibility.
\end{abstract}

\section{Introduction}
\label{sec:intro}

LLMs are increasingly deployed as \emph{controllers}: agents that
maintain context, plan over horizons, and produce actions in
response to a stream of observations.  This usage invites a
question that has a sharp, classical answer for linear feedback
systems but no settled answer for neural networks: \emph{what must
a controller contain in order to track an exogenous reference?}
For LTI systems, \citet{francis1976internal} answered with the
Internal Model Principle (IMP):  any controller that achieves
perfect asymptotic tracking of a class of disturbance or reference
signals must reproduce, internally, the dynamics that generate
those signals.  The IMP is an architectural necessity; it is
silent about \emph{representation}.  A model can live in a
state-space realization, in network weights, in activations, or in
tokens.

A widespread practitioner translation of the IMP to LLM agents
asserts that the representation is tokens: \emph{ask the model to
write down its hypothesis about the environment before acting, and
performance will improve}.  This translation is operational --- it
is implemented by editing a system prompt --- and it is the
implicit logic behind a substantial subset of agentic-system
design templates as well as the well-studied chain-of-thought (CoT)
prompting paradigm
\citep{wei2022chain,kojima2022zero,wang2023self}.  The
expectation is straightforward: articulating an internal model
should help, increasingly so with scale, and most of all on tasks
that genuinely require model-based prediction.

This paper tests that expectation directly and reports a
negative result.  We construct a clean sequence-prediction game in
which a hidden generator $f \colon \mathbb{N} \to \mathbb{Z}$
produces an integer target at each step $k$, the LLM observes the
history, emits a prediction $\hat y_k$, and the true $y_k$ is
appended.  We compare three matched system prompts: a
\reactive\ baseline that asks for one integer with no scaffolding,
an \imhyp\ prompt that explicitly invokes the IMP and instructs the
model to state a one-line hypothesis about $f(k)$ before predicting,
and an \imcot\ zero-shot CoT variant.

We then run two benchmarks against the same three prompts on four
open-weight model scales spanning roughly an order of magnitude.
A \emph{discovery} benchmark, with $30$ hand-curated tasks,
single-seed greedy decoding, and a paired-$t$ test over $10$ tasks
per difficulty band, is used to identify candidate effects.
A \emph{confirmation} benchmark, with $260$ programmatically
generated tasks across $30$ parametric families, multi-seed sampled
decoding at temperature $0.7$, and cluster-bootstrap confidence
intervals that resample tasks (the dominant variance cluster), is
used to test whether the candidate effects survive at the level of
statistical strength a careful reviewer requires.  Together the two
benchmarks comprise $5{,}760$ prompted trajectories.

The results, summarized in Sections~\ref{sec:results-v1}
and~\ref{sec:results-v2}, support three claims:

\begin{enumerate}[leftmargin=2em,itemsep=2pt,topsep=2pt]
\item \emph{Articulating an internal model never significantly
  helps.}  Across all $36$ \emph{difficulty $\times$ scale $\times$
  phrasing} cells in the two benchmarks, no paired test rejects
  the null in the helping direction at $\alpha = 0.05$.

\item \emph{The harm is robust at every scale.}  In the
  confirmation benchmark, the \imhyp\ prompt significantly
  degrades performance on $22$ of $24$ cells (paired $t$-tests on
  log-error, $p < 10^{-90}$ at the smallest scale and $p < 10^{-6}$
  at the largest).  Cluster-bootstrap CIs that resample tasks
  exclude zero on the medium-difficulty band at $7$B, $14$B and
  $32$B.

\item \emph{Chain-of-Thought articulation has an inverted-$U$ harm
  profile in scale.}  The mean log-harm of \imcot\ on hard tasks
  is $+0.02, +1.38, +1.60, +0.64$ at $7$B, $14$B, $32$B, $70$B
  respectively, peaking in the middle.  We propose, and discuss,
  a tokenized-commitment mechanism that predicts this shape:  a
  too-small model writes inert reasoning that the integer-emission
  step ignores, an intermediate model writes
  \emph{competent-looking} but wrong reasoning that the
  integer-emission step computes from, and a largest-tested model
  partially recovers by recognising and correcting the chain.
\end{enumerate}

A representative single-cell failure illustrates the mechanism:  on
the Fibonacci task at Qwen-7B, the \reactive\ condition attains zero
asymptotic error while the \imhyp\ condition attains an asymptotic
error of $71.17$.  Inspection of the raw replies reveals that the
\imhyp\ replies emit an incorrect recurrence on the
\verb|HYPOTHESIS:| line --- typically a wrong linear extrapolation
or a wrong arithmetic recurrence --- and then deterministically
compute integer predictions forward from the wrong rule, accumulating
error linearly in $k$.  The \reactive\ condition, by contrast, emits
the correct integer through pure pattern completion.  The same
phenomenon recurs at every scale.

We are careful to be precise about what these results refute and
what they do not.  They refute the prompting reading of the IMP --
the operational claim that asking an LLM to externalise an internal
model improves prediction.  They do not refute the IMP itself,
which is a theorem about representations available to a feedback
controller, not about prompts.  Section~\ref{sec:imp} works through
this distinction in detail.  Sections~\ref{sec:limitations}
and~\ref{sec:future} state the experimental scope honestly and
list pre-registered extensions.

\paragraph{Contributions.}
(i) A principled translation of the IMP to LLM-agent prompting,
together with three reasons it is not entitled to the strength of
the original theorem.
(ii) A reproducible $30$-task discovery benchmark and a
$260$-task confirmation benchmark with $30$ parametric task
families, deterministic given a sampling seed.
(iii) A multi-seed sampled-decoding protocol with cluster-bootstrap
CIs that resample tasks, addressing the standard "small-$n$
sensitivity" reviewer concern.
(iv) A robust, multi-scale negative result with an explicit
tokenized-commitment mechanism, including a previously unreported
inverted-$U$ harm pattern for CoT in scale.
(v) Code, prompts, and analysis pipeline released at
\url{https://github.com/Minds-R-Lab/phase_margin_llm}.

\section{The Internal Model Principle and Its Prompting Analogue}
\label{sec:imp}

We restate the classical IMP precisely so the prompting analogue
remains under tight scrutiny.

\paragraph{The classical principle.}  Consider an LTI plant $P$ and
controller $K$ in standard unity feedback, with reference $r$,
output $y$, and tracking error $e = r - y$.  Suppose $r$ belongs to
a class generated by an exosystem with state-space realization
$(A_e, B_e, C_e)$ -- e.g., $A_e$ has poles at $\{0, \pm j\omega_1,
\pm j\omega_2, \dots\}$ for steps and persistent sinusoids.
\citet[Theorem~1]{francis1976internal} establishes:
\emph{If the closed loop is internally stable and asymptotically
rejects the disturbance class generated by $(A_e, B_e, C_e)$, then
the dynamic modes of the exosystem must be reproduced inside the
loop.}  This is a necessity result -- internal stability plus
asymptotic tracking implies the internal model.

\paragraph{The prompting analogue.}  Translated to LLM agents, the
analogue reads:
\begin{quote}
\emph{If the agent must achieve low asymptotic error in tracking a
class of structured signals, the agent must, in some
representation, contain a model of those signals; therefore,
prompting the agent to externalize and use that model in tokens
should help.}
\end{quote}

Three observations show this analogue is not entitled to the
strength of the original.

\textbf{Modes, not language.}  The classical theorem is about which
poles appear in the loop in a fixed representation (state-space,
Laplace).  Translating "contains the model" to "writes the model
in words" requires that the textual articulation be representation-
preserving.  This is a strong assumption.  An LLM's implicit model
of integer sequences lives in network weights and activations;
articulating that implicit model in tokens may be lossy or
distortive.

\textbf{Asymptotic exact tracking, not approximate prediction.}  The
classical result is sharp at zero steady-state error.  An agent
doing finite-horizon prediction with average error as the loss
function is not in the regime where the sharp condition is
operative, and does not inherit the necessity-direction conclusion
\emph{a priori}.

\textbf{``Contains'' is not ``commits to''.}  The signal class
admits multiple realisations of its internal model.  A classical
controller may carry several modes simultaneously and combine them
linearly.  An LLM that has just emitted
\verb|HYPOTHESIS: f(k) = 5k + 1| has, in the next-token sense,
\emph{committed} to one realisation:  subsequent integer
predictions will be conditioned on that text and the autoregressive
decoding will lock in the implied computation.

These caveats predict that the prompting analogue is fragile.
The empirical question -- whether the analogue nevertheless yields
a useful heuristic -- is what the rest of this paper answers.

\section{Related Work}
\label{sec:related}

\paragraph{Chain-of-thought prompting.}
\citet{wei2022chain} introduced chain-of-thought (CoT) prompting and
demonstrated that prompting LLMs to produce intermediate reasoning
steps before answering improves accuracy on arithmetic, commonsense,
and symbolic-reasoning tasks, especially for the largest models in
their original setup ($\geq$ 100B parameters).
\citet{kojima2022zero} simplified this to a zero-shot
``Let's think step by step'' prompt, the exact phrasing we use in
our \imcot\ condition.  \citet{wang2023self} showed self-consistency
over multiple sampled chains improves CoT accuracy further.  More
recently, the meta-analysis of \citet{sprague2024cot} surveyed
$100+$ CoT papers and ran controlled evaluations on $20$ datasets;
their headline finding is that CoT helps mainly on math and
symbolic-execution tasks, with much smaller gains elsewhere.
Our paradigm -- short, deterministic integer-sequence prediction --
is symbolic in nature, the regime where CoT is supposed to help.
That \emph{both} the IMP-flavoured \imhyp\ and the
``Let's-think-step-by-step'' \imcot\ prompts \emph{hurt} on this
paradigm sharpens the negative-effects-of-CoT picture.

\paragraph{Multi-agent LLM failures.}
At a system level, \citet{cemri2025why} introduce MAST -- a taxonomy
of $14$ failure modes across $7$ multi-agent frameworks built from
expert annotation of $1{,}600+$ traces.  Their Category~I (system
design issues) includes failures driven by prompt structure and
role specification.  Our finding contributes a single-agent
specialisation: a system prompt that asks for an articulated
internal model can drive performance down on a task family that the
same model would solve with no scaffolding.

\paragraph{Control-theoretic perspectives on neural systems.}
A growing line of work imports control-theoretic concepts --
stability, observability, feedback -- into the analysis of neural
networks and LLM agents.  We take the opposite stance from much of
this work:  rather than borrowing a guarantee and asserting the
LLM inherits it, we test, empirically and conservatively, whether
one of the most foundational such guarantees translates at the
prompting level.  It does not.  This contributes an instructive
data point about which control-theoretic results have meaningful
prompt-level analogues.

\paragraph{Negative results with scale.}
Recent literature has documented several scale-dependent effects
in LLMs that contradict naive scaling intuitions, including
inverse-scaling failures (where larger models do worse at certain
tasks) and U-shaped scaling.  Our inverted-$U$ harm pattern for
CoT (Section~\ref{sec:cot-inverted-u}) is consistent with that
broader picture: the harm peaks where the model is competent
enough to write articulate but wrong reasoning, but not yet
competent enough to recognise and correct it.

\section{Benchmark Design and Methodology}
\label{sec:method}

We describe two benchmarks.  Benchmark v1 is a small,
hand-curated discovery benchmark; benchmark v2 is a large,
programmatically generated, multi-seed confirmation benchmark.
The same three system prompts and the same evaluation metric are
used in both benchmarks; only the task pool, decoding
configuration, and statistical analysis differ.

\subsection{System prompts (verbatim)}
\label{sec:prompts}

The central finding of this paper depends on subtle differences
between three system prompts.  We give them verbatim below; they
are identical across both benchmarks.

\paragraph{\reactive\ (no-scaffold baseline).}
\begin{Verbatim}[fontsize=\small,frame=single]
You are a sequence-prediction agent.  At each step you will be
shown all observations so far in the form
`k=0: y=...; k=1: y=...; ...` and asked to predict the next value
y_k.  Reply with EXACTLY one integer on a line by itself.  No
commentary, no labels, no quotes.
\end{Verbatim}

\paragraph{\imhyp\ (the canonical IMP-as-prompt).}
\begin{Verbatim}[fontsize=\small,frame=single]
You are a sequence-prediction agent operating under the Internal
Model Principle.  At each step you will be shown all observations
so far and asked to predict the next value y_k.  You MUST first
state a one-line HYPOTHESIS about the function f(k) that generated
the observed values, then on a NEW LINE write your single integer
prediction for y_k.

Reply format (exactly two lines):
    HYPOTHESIS: <your one-line hypothesis>
    <integer prediction>
\end{Verbatim}

\paragraph{\imcot\ (zero-shot chain-of-thought).}
\begin{Verbatim}[fontsize=\small,frame=single]
You are a sequence-prediction agent.  At each step you will be
shown all observations so far and asked to predict the next value
y_k.  Let's think step by step: briefly reason about what rule
could be generating these values and why your candidate next
value is consistent with it.  After your reasoning, write your
single integer prediction for y_k as the LAST line of your reply,
by itself, with no labels or punctuation.  The integer on the
last line will be parsed automatically.
\end{Verbatim}

The user prompt is identical across the three conditions:
\begin{Verbatim}[fontsize=\small,frame=single]
Observations so far: k=0: y=...; k=1: y=...; ...

Predict y_k.
\end{Verbatim}

\noindent
\reactive\ uses \verb|max_tokens = 96|; \imhyp\ uses
\verb|max_tokens = 96|; \imcot\ uses \verb|max_tokens = 256|
(reasoning is longer).  This is the only systematic difference in
decoding configuration between conditions and we revisit it in
Section~\ref{sec:limitations}.

\subsection{Common evaluation metric}
\label{sec:metric}

Each cell consists of a $K$-step trajectory.  At step $k$, the
model is shown the history of true observations $\{(j, y_j)\}_{j<k}$
and asked to predict $y_k$.  The history is then advanced by
appending the \emph{true} $y_k$ (not the prediction), so the loop
is open in the prediction direction.  The asymptotic error of the
cell is the mean absolute error over the last $\textsc{asym}$
steps:
\begin{equation}
  \errk \;=\; \frac{1}{\textsc{asym}}
  \sum_{k=K-\textsc{asym}}^{K-1} \big| \hat y_k - y_k \big|.
\end{equation}
For statistical comparisons we work with
$\log(\errk + \varepsilon)$ with $\varepsilon = 0.5$ to keep the
metric defined when one or both conditions score zero.

\subsection{Benchmark v1: $30$-task discovery benchmark}
\label{sec:method-v1}

The discovery benchmark contains $30$ deterministic, integer-valued
generators, balanced across three difficulty bands of $10$ tasks
each.  Difficulty was hand-assigned: \emph{easy} tasks are those
on which a no-reasoning baseline (echo recent observation) tends
to be approximately correct; \emph{medium} tasks require explicit
pattern recognition but only one symbolic step to compute the next
value; \emph{hard} tasks require non-trivial computation even given
a correct hypothesis.  Trajectory length is $K = 14$ with
$\textsc{asym} = 6$, except for one geometric task capped at $K = 12$
to avoid integer overflow in the prompt.  Decoding is greedy
($\text{temperature} = 0$, single seed).  The full task list and
the exact generators are given in Table~\ref{tab:tasks-v1} of
Appendix~\ref{app:tasks}.

The discovery benchmark uses paired $t$-tests on log-error per
(band, IM-condition) pair, with pairing over the $10$ tasks.
$n = 10$ per test.

\subsection{Benchmark v2: $260$-task confirmation benchmark}
\label{sec:method-v2}

A reviewer of the v1 benchmark observes, correctly, that
single-seed greedy decoding gives no within-cell variance and that
$n = 10$ tasks per band is sensitive to the specific tasks chosen.
Benchmark v2 is engineered to address both prongs of this critique
simultaneously.

\paragraph{Programmatic task families.}  Each difficulty band is
a pool drawn from $10$ parametric generator families.  For example,
the \emph{easy} band families are: a constant integer, a positive
linear function $a k + b$, a negative-slope linear function, a
period-$2$ cycle, a period-$3$ cycle, two period-$4$ cycle variants,
$k \bmod n$, an arithmetic progression $a k$, and an arithmetic
progression with offset.  The \emph{medium} families include
quadratic, quadratic with offset, negative quadratic, period-$5$
and period-$6$ cycles, piecewise linear, zigzag $k(-1)^k s$,
$\lfloor k/n \rfloor$, triangle wave, and linear-plus-period.  The
\emph{hard} families include Fibonacci with random initial
conditions, capped geometric $c \cdot 2^k$, $k(k-1)$, triangle
numbers, random digit cycles, $k^3 \bmod n$, conditional rules,
regime-switching cycles, sums of two coprime periods, and general
quadratics with mixed-sign coefficients.  Each family is
parameterized by integer ranges chosen to keep prompts bounded
(target absolute value capped at $50{,}000$) and to be
distinguishable from one another.

A task is generated by drawing a parameter tuple from a family.
The full task pool size is configurable; we use $100$ tasks per
band on the smallest model, scaling down to $20$ on the largest
to fit overnight wall-time.  The set is deterministic given a
sampling seed: the same \verb|--task-seed| produces the identical
task pool across machines.  All four model runs in this paper
use \verb|--task-seed=2026|.

\paragraph{Multi-seed sampled decoding.}  Each cell is run with
multiple LLM seeds at \verb|temperature = 0.7|.  The number of
seeds per cell is $5$ on the two smallest models, $3$ on the two
larger models, again to fit wall-time.  Pairing is over (task,
seed), so each (band, IM-condition) test has up to $500$ paired
observations rather than $10$.  Trajectory length is $K = 12$ with
$\textsc{asym} = 5$ throughout.

\paragraph{Statistical analysis.}  For each (scale, band,
IM-condition) triple we report:
\begin{itemize}[leftmargin=2em,itemsep=0pt,topsep=2pt]
\item a paired $t$-test on log-error over (task, seed),
\item a cluster-bootstrap $95\%$ CI on the median log-difference,
  resampling at the \emph{task} level (the dominant variance
  cluster), with $5{,}000$ bootstrap replicates.  Resampling tasks
  --- not seeds --- is what addresses the
  ``sensitive to which tasks?'' reviewer concern directly,
\item a within-task vs.\ across-task variance decomposition.
\end{itemize}

\paragraph{Implementation.}  The benchmark writes per-cell results
incrementally to a JSON-Lines file.  An interrupted overnight run
is resumed by re-invoking the same command; cells already present
in the JSONL are skipped.  Total wall-time across the four scales
on a single H100 is approximately $26$ hours.

\subsection{Models}
\label{sec:models}

We test four open-weight model scales served via Ollama on a
single NVIDIA H100~80GB:
\begin{itemize}[leftmargin=2em,itemsep=0pt,topsep=2pt]
\item \textbf{Qwen2.5-7B-Instruct} ($7.6$B parameters)
\item \textbf{Qwen2.5-14B-Instruct} ($14.7$B parameters)
\item \textbf{Qwen2.5-32B-Instruct} ($32.5$B parameters)
\item \textbf{Llama3.1-70B-Instruct, q4\_K\_M} ($70$B parameters,
  $4$-bit quantized)
\end{itemize}
The four scales span roughly an order of magnitude in parameter
count.  All four are open-weight; the $70$B is the only quantized
model in the set.  Frontier-scale closed models (Claude, GPT,
Gemini) are not tested here; see Section~\ref{sec:future}.

\section{Pre-registered Hypotheses}
\label{sec:hypotheses}

The naive prompting reading of the IMP, combined with the
chain-of-thought literature, suggests three hypotheses.  We
state them before reporting results.

\textbf{H1.}  IMP prompts hurt small models on tasks they can
already solve reactively, but the harm decays monotonically with
scale.

\textbf{H2.}  IMP prompts help small models on hard tasks
(replicating the chain-of-thought literature for math/symbolic
reasoning).

\textbf{H3.}  The joint behaviour shows a significant
\emph{difficulty $\times$ condition $\times$ scale} interaction.

The empirical findings, reported next, refine H1 (the harm decays
in some bands and \emph{migrates} in others), refute H2 (no help
emerges anywhere, at any scale), and confirm a modified form of
H3 (the interaction is present but in the
harm-migration direction, and a previously unreported inverted-$U$
appears for CoT).

\section{Results: Discovery Benchmark}
\label{sec:results-v1}

\subsection{Main table}
Table~\ref{tab:main-v1} reports the mean asymptotic error per
(scale, difficulty, condition) cell on the v1 benchmark.

\begin{table}[t]
\centering
\caption{Discovery benchmark (v1, $n = 10$ tasks per cell, single
seed, greedy decoding).  Mean asymptotic error per
(scale, difficulty, condition).  Lower is better.  Bold marks
cells where a paired $t$-test on log-error against the same row's
\reactive\ baseline is significant at $p < 0.05$ in the
\emph{hurts} direction.  No cell shows the help direction at any
significance level.}
\label{tab:main-v1}
\small
\begin{tabular}{ll rrr}
\toprule
Model & Difficulty & \reactive & \imhyp & \imcot \\
\midrule
\multirow{3}{*}{Qwen2.5-7B}
  & easy   & $0.15$ & $\mathbf{6.08}^{\hurts}$  & $0.57$ \\
  & medium & $1.58$ & $\mathbf{5.22}^{\hurts}$  & $1.90$ \\
  & hard   & $5.58$ & $30.03$                   & $7.05$ \\
\addlinespace
\multirow{3}{*}{Qwen2.5-14B}
  & easy   & $0.00$ & $\mathbf{1.72}^{\hurts}$  & $\mathbf{1.32}^{\hurts}$ \\
  & medium & $0.43$ & $\mathbf{9.22}^{\hurts}$  & $\mathbf{2.57}^{\hurts}$ \\
  & hard   & $5.40$ & $5.43$                    & $45.12$ \\
\addlinespace
\multirow{3}{*}{Qwen2.5-32B}
  & easy   & $0.00$ & $0.75$                    & $0.25$ \\
  & medium & $1.73$ & $\mathbf{2.48}^{\hurts}$  & $2.30$ \\
  & hard   & $5.57$ & $20.85$                   & $\mathbf{13.53}^{\hurts}$ \\
\addlinespace
\multirow{3}{*}{Llama3.1-70B-q4}
  & easy   & $0.00$ & $0.00$\,\textsuperscript{(tied)}  & $0.02$ \\
  & medium & $0.37$ & $\mathbf{0.97}^{\hurts}$  & $\mathbf{1.83}^{\hurts}$ \\
  & hard   & $6.77$ & $8.03$                    & $11.25$ \\
\bottomrule
\end{tabular}
\end{table}

The discovery benchmark identifies the headline pattern: at every
scale, on the medium difficulty band, \imhyp\ significantly hurts
($p$ values $0.001, 0.010, 0.043, 0.025$ across the four scales,
combining via Stouffer's method to $Z = 5.06$, $p < 10^{-6}$).
On the easy band, the harm is significant at the smaller scales
and ties at the floor for the largest.  On the hard band, the v1
benchmark is underpowered; significant effects appear only on the
$32$B-CoT cell.

\subsection{Cross-scale trend on the easy band}
The mean log-error gap of \imhyp\ vs.\ \reactive\ on the easy band
shrinks $+1.41 \to +1.08 \to +0.51 \to 0.00$ across the four
scales.  Treating the four points as a function of
$\log(\text{parameter count})$ gives Spearman $\rho = -1.00$ (perfect
monotone decay) and Pearson $r = -0.998$ ($p = 0.002$ two-sided).
The medium band gives a directionally consistent but
not-quite-significant decay (Spearman $-0.80$, Pearson $-0.92$).
The pattern is interpretable as the easy-band reactive baseline
saturating at the floor of zero error, leaving no room for the
prompt to harm; it does not generalize cleanly to the medium band,
where reactive performance still has slack at every scale.

\subsection{Catastrophic single-cell failures}
A subset of cells exhibit a particularly informative failure mode:
\reactive\ achieves zero asymptotic error and the articulated
condition produces a persistent error of tens to hundreds of units.
Such cells rule out the explanation ``the model cannot do the
task.''  Table~\ref{tab:catas-v1} lists the most striking cases at
each scale.  A representative example: on Fibonacci with Qwen-7B,
\reactive\ scores $0.00$ while \imhyp\ scores $71.17$.  Inspection
of the raw replies shows the model emits a wrong recurrence on the
\verb|HYPOTHESIS:| line and computes integer predictions
deterministically forward from it.  We return to this mechanism in
Section~\ref{sec:mechanism}.

\begin{table}[t]
\centering
\caption{Catastrophic single-cell failures (v1).  Cells where
\reactive\ scores at floor and the articulated condition does not.
The phenomenon recurs at every scale, including the $32$B model.}
\label{tab:catas-v1}
\small
\begin{tabular}{ll rrr}
\toprule
Model & Task & \reactive & \imhyp & \imcot \\
\midrule
Qwen2.5-7B  & \texttt{e04\_linear\_5k}        & $0.00$ & $41.67$  & $0.00$  \\
Qwen2.5-7B  & \texttt{h01\_fibonacci}         & $0.00$ & $71.17$  & $0.00$  \\
Qwen2.5-7B  & \texttt{h02\_geometric\_2}      & $0.00$ & $160.00$ & $0.00$  \\
Qwen2.5-14B & \texttt{m02\_quadratic\_offset} & $0.00$ & $66.17$  & $7.83$  \\
Qwen2.5-14B & \texttt{h02\_geometric\_2}      & $0.00$ & $0.00$   & $384.00$ \\
Qwen2.5-32B & \texttt{m02\_quadratic\_offset} & $0.00$ & $8.83$   & $16.50$ \\
Qwen2.5-32B & \texttt{h02\_geometric\_2}      & $0.00$ & $159.50$ & $0.00$  \\
Qwen2.5-32B & \texttt{h03\_squares\_minus\_k} & $0.00$ & $0.00$   & $68.67$ \\
\bottomrule
\end{tabular}
\end{table}

\section{Results: Confirmation Benchmark}
\label{sec:results-v2}

\subsection{Main table}

Benchmark v2 increases $n$ per cell from $10$ to between $60$ and
$500$ paired observations, replaces hand-curated tasks with
programmatically sampled tasks, and introduces multi-seed sampled
decoding.  The mean asymptotic errors are reported in
Table~\ref{tab:main-v2}.

\begin{table}[t]
\centering
\caption{Confirmation benchmark (v2, programmatic tasks,
multi-seed sampled decoding at $T = 0.7$).  Mean asymptotic error
per (scale, difficulty, condition).  Sample size per condition,
\textsc{n}\textsubscript{cells}, is given in the second column.
Across this corpus, $22$ of $24$ paired tests reject the null in
the \emph{hurts} direction at $\alpha = 0.05$; zero reject in the
\emph{helps} direction.  Bold marks $p < 0.05$ hurt significance
of the paired $t$-test.}
\label{tab:main-v2}
\small
\setlength{\tabcolsep}{4pt}
\begin{tabular}{ll r rrr}
\toprule
Model & Difficulty & \textsc{n}\textsubscript{cells}
  & \reactive & \imhyp & \imcot \\
\midrule
\multirow{3}{*}{Qwen2.5-7B}
  & easy   & 500 & $\sim 0.08$ & $\mathbf{2.6}^{\hurts}$  & $\mathbf{0.95}^{\hurts}$ \\
  & medium & 500 & $1.83$      & $\mathbf{6.67}^{\hurts}$ & $\mathbf{1.94}^{\hurts}$ \\
  & hard   & 500 & $7.56$      & $\mathbf{72.93}^{\hurts}$ & $8.82$ \\
\addlinespace
\multirow{3}{*}{Qwen2.5-14B}
  & easy   & 400 & $0.00$      & $\mathbf{1.60}^{\hurts}$  & $\mathbf{1.23}^{\hurts}$ \\
  & medium & 400 & $0.19$      & $\mathbf{4.12}^{\hurts}$  & $\mathbf{2.52}^{\hurts}$ \\
  & hard   & 400 & $6.12$      & $\mathbf{11.90}^{\hurts}$ & $\mathbf{21.78}^{\hurts}$ \\
\addlinespace
\multirow{3}{*}{Qwen2.5-32B}
  & easy   & 180 & $0.00$      & $\mathbf{1.30}^{\hurts}$  & $\mathbf{0.46}^{\hurts}$ \\
  & medium & 180 & $0.33$      & $\mathbf{4.63}^{\hurts}$  & $\mathbf{6.34}^{\hurts}$ \\
  & hard   & 180 & $47.16$     & $\mathbf{45.80}^{\hurts}$ & $\mathbf{29.24}^{\hurts}$ \\
\addlinespace
\multirow{3}{*}{Llama3.1-70B-q4}
  & easy   &  60 & $0.07$      & $\mathbf{0.19}^{\hurts}$  & $\mathbf{1.12}^{\hurts}$ \\
  & medium &  60 & $0.40$      & $\mathbf{1.32}^{\hurts}$  & $\mathbf{2.14}^{\hurts}$ \\
  & hard   &  60 & $6.79$      & $10.64$                   & $\mathbf{17.67}^{\hurts}$ \\
\bottomrule
\end{tabular}
\end{table}

\subsection{Paired $t$-tests}

The paired $t$-statistics on log-error per cell, against the
\reactive\ baseline, are reported in Table~\ref{tab:ttests-v2}.
Recall:  positive $t$ here means \imhyp\ or \imcot\ has
\emph{larger} log-error than \reactive, i.e., the IM condition
hurts.

\begin{table}[t]
\centering
\caption{Confirmation benchmark paired $t$-tests on log-error
$\log(\errk + 0.5)$.  \reactive\ vs.\ each IM condition, by scale
and difficulty.  Positive $t$ means the IM condition hurts.  At
$\alpha = 0.05$, $22$ of $24$ tests reject the null in the
\emph{hurts} direction; the only non-significant cells are the
$7$B-hard-CoT and $70$B-hard-Hyp.}
\label{tab:ttests-v2}
\small
\setlength{\tabcolsep}{4pt}
\begin{tabular}{l l rr rr}
\toprule
Model & Difficulty
  & \multicolumn{2}{c}{vs.\ \imhyp}
  & \multicolumn{2}{c}{vs.\ \imcot} \\
\cmidrule(lr){3-4}\cmidrule(lr){5-6}
 & & $t$ & $p$ & $t$ & $p$ \\
\midrule
\multirow{3}{*}{Qwen2.5-7B}
  & easy   & $+27.0$ & $\mathbf{8.7\!\cdot\!10^{-100}}^{\hurts}$ & $+6.16$ & $\mathbf{1.5\!\cdot\!10^{-9}}^{\hurts}$ \\
  & medium & $+23.2$ & $\mathbf{4.4\!\cdot\!10^{-81}}^{\hurts}$  & $+3.30$ & $\mathbf{0.0010}^{\hurts}$ \\
  & hard   & $+16.7$ & $\mathbf{7.2\!\cdot\!10^{-50}}^{\hurts}$  & $+0.50$ & $0.62$ \\
\multirow{3}{*}{Qwen2.5-14B}
  & easy   & $+16.6$ & $\mathbf{2.3\!\cdot\!10^{-47}}^{\hurts}$  & $+15.4$ & $\mathbf{3.2\!\cdot\!10^{-42}}^{\hurts}$ \\
  & medium & $+27.2$ & $\mathbf{5.6\!\cdot\!10^{-93}}^{\hurts}$  & $+19.8$ & $\mathbf{2.3\!\cdot\!10^{-61}}^{\hurts}$ \\
  & hard   & $+12.8$ & $\mathbf{1.3\!\cdot\!10^{-31}}^{\hurts}$  & $+15.6$ & $\mathbf{5.0\!\cdot\!10^{-43}}^{\hurts}$ \\
\multirow{3}{*}{Qwen2.5-32B}
  & easy   & $+11.1$ & $\mathbf{3.1\!\cdot\!10^{-22}}^{\hurts}$  & $+6.57$ & $\mathbf{5.3\!\cdot\!10^{-10}}^{\hurts}$ \\
  & medium & $+12.3$ & $\mathbf{1.8\!\cdot\!10^{-25}}^{\hurts}$  & $+9.29$ & $\mathbf{5.2\!\cdot\!10^{-17}}^{\hurts}$ \\
  & hard   & $+6.33$ & $\mathbf{1.9\!\cdot\!10^{-9}}^{\hurts}$   & $+9.95$ & $\mathbf{7.9\!\cdot\!10^{-19}}^{\hurts}$ \\
\multirow{3}{*}{Llama3.1-70B-q4}
  & easy   & $+2.67$ & $\mathbf{0.0098}^{\hurts}$ & $+3.35$ & $\mathbf{0.0014}^{\hurts}$ \\
  & medium & $+5.61$ & $\mathbf{5.7\!\cdot\!10^{-7}}^{\hurts}$  & $+6.28$ & $\mathbf{4.5\!\cdot\!10^{-8}}^{\hurts}$ \\
  & hard   & $+1.61$ & $0.11$                    & $+3.15$ & $\mathbf{0.0026}^{\hurts}$ \\
\bottomrule
\end{tabular}
\end{table}

\subsection{Cluster-bootstrap confidence intervals}

A paired $t$-test pools across tasks; an alert reviewer asks how
sensitive the conclusion is to the specific tasks sampled.
Table~\ref{tab:bootstrap-v2} reports cluster-bootstrap $95\%$
confidence intervals on the median log-difference, resampling at
the task level with $5{,}000$ replicates.  An interval that
excludes zero is robust to the specific task sample.

\begin{table}[t]
\centering
\caption{Cluster-bootstrap $95\%$ CI on the median
log-difference $\log(\errk^{\text{IM}} + 0.5)
- \log(\errk^{\reactive} + 0.5)$, resampling at the task
level ($5{,}000$ replicates).  An interval excluding zero in the
positive direction implies the IM condition robustly hurts even
under task resampling.  The bootstrap is markedly more
conservative than the paired $t$-test because most tasks have
median log-difference exactly zero (both conditions tied);
the test detects the heavy tail of large-error tasks while the
median-of-bootstrap stays at zero unless the tail is large
relative to the task pool.}
\label{tab:bootstrap-v2}
\small
\setlength{\tabcolsep}{3pt}
\begin{tabular}{l l ll}
\toprule
Model & Difficulty
  & vs.\ \imhyp\ \,(median, $95\%$~CI)
  & vs.\ \imcot\ \,(median, $95\%$~CI) \\
\midrule
\multirow{3}{*}{Qwen2.5-7B}
  & easy   & $\mathbf{1.44}\ [+1.10,\,+1.69]^{\hurts}$ & $0.00\ [\,0,\,0\,]$ \\
  & medium & $\mathbf{1.00}\ [+0.76,\,+1.29]^{\hurts}$ & $0.00\ [\,0,\,0\,]$ \\
  & hard   & $\mathbf{0.71}\ [+0.40,\,+1.31]^{\hurts}$ & $0.00\ [\,0,\,0\,]$ \\
\multirow{3}{*}{Qwen2.5-14B}
  & easy   & $0.34\ [\,0,\,+0.96\,]$                   & $0.00\ [\,0,\,+0.59\,]$ \\
  & medium & $\mathbf{1.67}\ [+1.22,\,+1.95]^{\hurts}$ & $\mathbf{0.96}\ [+0.35,\,+1.25]^{\hurts}$ \\
  & hard   & $\mathbf{0.28}\ [+0.09,\,+0.41]^{\hurts}$ & $\mathbf{0.43}\ [+0.29,\,+0.95]^{\hurts}$ \\
\multirow{3}{*}{Qwen2.5-32B}
  & easy   & $0.17\ [\,0,\,+0.96\,]$                   & $0.00\ [\,0,\,0\,]$ \\
  & medium & $\mathbf{0.79}\ [+0.37,\,+1.44]^{\hurts}$ & $0.34\ [\,0,\,+1.22\,]$ \\
  & hard   & $\mathbf{0.32}\ [+0.08,\,+0.82]^{\hurts}$ & $\mathbf{0.46}\ [+0.11,\,+2.24]^{\hurts}$ \\
\multirow{3}{*}{Llama3.1-70B-q4}
  & easy   & $0.00\ [\,0,\,0\,]$                       & $0.00\ [\,0,\,0\,]$ \\
  & medium & $0.25\ [\,0,\,+0.62\,]$                   & $\mathbf{0.45}\ [+0.06,\,+1.19]^{\hurts}$ \\
  & hard   & $0.00\ [\,0,\,+0.18\,]$                   & $0.00\ [\,0,\,+0.18\,]$ \\
\bottomrule
\end{tabular}
\end{table}

The most robust headline that survives both the parametric
$t$-test and the non-parametric task-cluster bootstrap is:
\emph{the IM-Hypothesis prompt significantly hurts on the
medium-difficulty band at $7$B, $14$B, and $32$B}, with bootstrap
CIs $[+0.76,+1.29]$, $[+1.22,+1.95]$, and $[+0.37,+1.44]$
respectively.  At $70$B the bootstrap CI on medium-Hyp is
$[0,+0.62]$, just touching zero, although the paired $t$-test is
highly significant ($p < 10^{-6}$).

We emphasize the principled reading of the gap between the two
tests.  The paired $t$-test detects mean log-difference; the
bootstrap CI on the median is dominated by tasks where both
conditions tie.  For \imcot\ in particular, a large fraction of
tasks have identical error under both conditions (CoT often
reduces to reactive when the reasoning chain ends with the same
integer); the median-of-bootstrap is therefore zero while the
mean is positive.  This is a feature of the task pool: harm is
concentrated on a subset of tasks, not uniformly spread.

\subsection{Inverted-$U$ harm profile of CoT in scale}
\label{sec:cot-inverted-u}

Reading the \imcot\ column of Table~\ref{tab:main-v2} by scale on
the hard band reveals a non-monotone pattern:  mean log-harm
$+0.02 \to +1.38 \to +1.60 \to +0.64$ at $7$B, $14$B, $32$B,
$70$B respectively, peaking in the middle.  The same shape appears
on the medium band, less pronounced.  Figure~\ref{fig:invertedU}
plots the four points by scale on the hard and medium bands.
The shape is consistent with a tokenized-commitment account:

\begin{itemize}[leftmargin=2em,itemsep=0pt,topsep=2pt]
\item The smallest model ($7$B) often fails to follow the CoT
  scaffolding tightly; its reasoning chain is short, generic, and
  ends in an integer largely independent of the chain's content.
  In effect, \imcot\ collapses to \reactive\ for the smallest
  model, and the harm is small.

\item The intermediate models ($14$B, $32$B) write
  \emph{competent-looking} reasoning chains -- recognisably
  sequence-pattern-fitting prose -- that often arrive at a wrong
  conclusion and pull the integer prediction with them.  The
  harm is large.

\item The largest model ($70$B-q4) writes longer, more careful
  reasoning chains that more often catch their own errors and
  re-converge to the correct integer.  The harm is intermediate.
\end{itemize}

This profile is a positive empirical contribution that the v1
benchmark, with its small $n$, could not see.

\begin{table}[t]
\centering
\caption{Inverted-$U$ profile of \imcot\ harm in scale, on the
hard and medium difficulty bands.  Reported values are the mean
log-difference of \imcot\ vs.\ \reactive\ on the v2 benchmark.
Both bands peak at intermediate scale ($14$B--$32$B) and drop
toward both extremes.}
\label{fig:invertedU}
\small
\begin{tabular}{l rrrr}
\toprule
Difficulty band & Qwen-7B & Qwen-14B & Qwen-32B & Llama-70B-q4 \\
\midrule
hard   & $+0.02$ & $+1.38$ & $+1.60$ & $+0.64$ \\
medium & $+0.07$ & $+1.03$ & $+1.10$ & $+0.75$ \\
\bottomrule
\end{tabular}
\end{table}

\subsection{Variance decomposition}

Within-task variance and across-task variance of log-error are
reported in Appendix Table~\ref{tab:variance-v2}.  The headline
fact:  the dominant variance source is across-task, not
within-task, by factors of $5$--$20\times$ on the \reactive\ and
\imhyp\ conditions.  This implies that pairing by task -- the v1
strategy -- captured the right signal, and that multi-seed
decoding contributes additional power but does not change the
qualitative conclusion.  An exception is \imcot, where within-task
variance grows substantially due to the diversity of reasoning
chains across seeds; this is reflected in the conservative
bootstrap CIs for the \imcot\ column of
Table~\ref{tab:bootstrap-v2}.

\section{Mechanism: Tokenized Commitment}
\label{sec:mechanism}

Three independent observations point to the same mechanism for
the harm.

\textbf{Observation 1.}  Catastrophic single-cell failures
(Tables~\ref{tab:catas-v1} and~\ref{tab:catas-v2}) all share the
same structural form:  the model emits a wrong rule on the
\verb|HYPOTHESIS:| line and then deterministically computes
integer predictions forward from the wrong rule.  The errors are
not noisy -- they are exact arithmetic consequences of the wrong
hypothesis text.  On Fibonacci with Qwen-7B, the recurring wrong
hypothesis is a linear extrapolation; on geometric-2 with Qwen-32B,
the recurring wrong hypothesis is a wrong base.  We provide raw
replies for two cells in Appendix~\ref{app:replies}.

\textbf{Observation 2.}  The catastrophic failures recur at every
scale tested in v1, including at $32$B.  This is inconsistent with
the simplest "small models can't follow IMP scaffolding"
explanation:  a $32$B model that demonstrates correct integer
prediction under \reactive\ on a task and incorrect prediction
under \imhyp\ on the same task is not failing for capability
reasons.

\textbf{Observation 3.}  The inverted-$U$ in scale for \imcot\
fits the commitment account:  CoT prompts only harm when the
articulation is detailed enough to commit the decoder to a
particular forward computation.  At the smallest scale,
articulation is sketchy and uncommitted; at the largest,
articulation is corrigible.  Maximum commitment occurs at
intermediate scale.

\paragraph{Tokenized-commitment account.}  Within an
autoregressive decoder, generated tokens are part of the prefix
on which subsequent tokens are conditioned.  The \verb|HYPOTHESIS:|
line is therefore a soft prefix on the integer-prediction line:
once a particular rule is articulated, the decoder's next-token
distribution over the integer line is conditioned on the rule.  If
the rule is correct, conditioning helps:  the integer is the
arithmetic image of the rule.  If the rule is wrong, conditioning
hurts:  the integer is the arithmetic image of the wrong rule, and
no further information helps recover from it.  Crucially,
\reactive\ never produces such a textual prefix.  Its prediction
is conditioned only on the history of integer pairs, not on a
self-articulated rule.  In effect, the implicit pattern-matcher
encoded in the network's weights and activations is more accurate
on this paradigm than the rule the model articulates as text.

In control-theoretic terms:  the IM-Hypothesis prompt does not
expose the implicit internal model the network already has; it
forces the decoder to \emph{generate} a separate, articulated
internal model in tokens, which then becomes the operative model
for downstream prediction.  In a fraction of cells, this articulated
model is worse than the implicit one --- and crucially, once
articulated, the wrong model is committed to.

\paragraph{What the mechanism predicts.}  This account makes
several testable predictions.  Most directly, sampling several
hypotheses per step and aggregating (self-consistency, in the
sense of \citealp{wang2023self}) should restore performance to
near \reactive.  Decomposing harm cell-by-cell should show that
cells where the articulated hypothesis is correct match or beat
\reactive, while cells where the hypothesis is wrong account for
the entire \imhyp$-$\reactive\ gap.  The self-consistency
prediction is pre-registered in Section~\ref{sec:future}; the
hypothesis-correctness decomposition is reported next.

\subsection{Direct test: hypothesis-correctness decomposition}
\label{sec:hyp-correctness}

We perform a free post-hoc analysis on the v2 cells.  For every
\imhyp\ cell we extract the \verb|HYPOTHESIS:|\ line at every
step, parse it with a regex grammar of $20+$ patterns covering
linear, constant, $k \bmod n$, $k^2$, $(k\pm a)^2$, general
quadratic, $k(k-1)$, triangle numbers $k(k+1)/2$, $k^3 \bmod n$,
geometric $c \cdot 2^k$, triangle wave, $\lfloor k/n \rfloor$,
zigzag, Fibonacci with explicit initial conditions, and several
period/cycle phrasings (bracket lists, ``alternates between $A$
and $B$'', ``repeats every $N$ terms with the pattern $A, B,
\dots$'').  Each parsed rule is evaluated at the same step and
compared to (a) the true target $y_k$ and (b) the model's own
integer prediction $\hat y_k$.  We bin each cell on two binary
axes over the asymptotic window:
\emph{rule-correct} (does the parsed rule predict $y_k$ correctly
on $\geq 66\%$ of asym-window steps?)\ and
\emph{rule-followed} (does the model's integer match the rule's
prediction on $\geq 66\%$ of asym-window steps?).
A cell where either rate falls in the borderline $(0.34, 0.66)$
region is binned ``mixed''; cells where the parser cannot read
any rule in the window are binned ``unparseable''.

Two failure modes are possible under the \imhyp\ prompt:
\begin{enumerate}[leftmargin=2em,itemsep=0pt,topsep=2pt]
\item \textbf{Wrong rule, committed} (\textit{wrong\_followed}):
  the model articulates an incorrect rule and emits the integer
  the rule predicts.  Predicted by the original
  tokenized-commitment account.
\item \textbf{Correct rule, not applied}
  (\textit{correct\_not\_followed}): the model articulates a
  correct rule but emits an integer that disagrees with it.  This
  decoupling between rule-articulation and rule-application has
  no analogue under \reactive\ and is not predicted by the
  original account.
\end{enumerate}
The grid is filled by \textit{correct\_followed}
(commitment-faithful good case, expected to match \reactive) and
\textit{wrong\_not\_followed} (residual category).

Table~\ref{tab:bin4way} reports the $2{\times}2$ grid for all
four scales.

\begin{table}[t]
\centering
\caption{Four-way (rule-correct $\times$ rule-followed)
decomposition on v2 \imhyp\ cells.  Reported per (scale, bin):
sample size $n$, mean asymptotic error of the IM cell, mean
asymptotic error of the same task's reactive cell (paired by
task), and the gap.  Bins use a $66\%$-of-asym-window threshold
on each axis; ``mixed'' is borderline on either axis;
``unparseable'' denotes cells where the regex grammar could not
read a rule.}
\label{tab:bin4way}
\small
\setlength{\tabcolsep}{4pt}
\begin{tabular}{l l r rr r}
\toprule
Model & Bin & $n$ & mean IM & mean \reactive & gap \\
\midrule
\multirow{6}{*}{Qwen2.5-7B}
  & correct\_followed       &  34 & $106.51$ & $0.25$ & $+106.26$ \\
  & correct\_not\_followed  &  95 & $11.34$  & $0.21$ & $+11.13$ \\
  & wrong\_followed         &  25 & $218.00$ & $0.15$ & $+217.85$ \\
  & wrong\_not\_followed    & 129 & $22.33$  & $5.21$ & $+17.12$ \\
  & mixed                   & 139 & $8.02$   & $0.55$ & $+7.47$ \\
  & unparseable             &1078 & $26.53$  & $3.74$ & $+22.79$ \\
\addlinespace
\multirow{6}{*}{Qwen2.5-14B}
  & correct\_followed       &  81 & $0.91$   & $0.03$ & $+0.88$ \\
  & correct\_not\_followed  &  48 & $3.37$   & $0.10$ & $\mathbf{+3.27}$ \\
  & wrong\_followed         &   9 & $2.93$   & $0.91$ & $+2.03$ \\
  & wrong\_not\_followed    & 132 & $3.21$   & $1.05$ & $+2.16$ \\
  & mixed                   & 105 & $3.02$   & $0.22$ & $+2.80$ \\
  & unparseable             & 825 & $7.33$   & $2.85$ & $+4.48$ \\
\addlinespace
\multirow{6}{*}{Qwen2.5-32B}
  & correct\_followed       &  87 & $0.46$   & $0.02$ & $+0.45$ \\
  & correct\_not\_followed  &   3 & $1.87$   & $1.29$ & $+0.58$ \\
  & wrong\_followed         &   6 & $3.97$   & $2.07$ & $+1.90$ \\
  & wrong\_not\_followed    &  53 & $2.19$   & $17.02$& $-14.82$\,$^{\star}$ \\
  & mixed                   &  58 & $6.52$   & $0.71$ & $+5.81$ \\
  & unparseable             & 333 & $26.27$  & $22.79$& $+3.48$ \\
\addlinespace
\multirow{6}{*}{Llama3.1-70B-q4}
  & correct\_followed       &  22 & $0.082$  & $0.045$& $\mathbf{+0.04}$ \\
  & correct\_not\_followed  &   1 & $0.60$   & $0.73$ & $-0.13$ \\
  & wrong\_followed         &   2 & $4.00$   & $2.63$ & $+1.37$ \\
  & wrong\_not\_followed    &  26 & $1.31$   & $0.85$ & $+0.45$ \\
  & mixed                   &  14 & $0.56$   & $0.61$ & $-0.05$ \\
  & unparseable             & 115 & $5.89$   & $3.46$ & $+2.43$ \\
\bottomrule
\end{tabular}
\par\smallskip
\footnotesize
$^{\star}$ Negative gap is a within-task selection artefact on the
heavy-tailed hard band (\texttt{h07\_cubic\_mod100} has reactive
error $\sim 50$); it is not evidence that \imhyp\ helps.
\end{table}

The reading by scale is informative.

\textbf{Llama-70B-q4: the original tokenized-commitment story
holds sharply.}  When the parser confirms a correct rule
\emph{and} the model emits an integer consistent with it, the
asymptotic error is at the reactive floor: $0.082$ for IM vs.\
$0.045$ for reactive, gap $+0.04$.  Wrong-rule-followed cells
(small $n$, but directionally consistent) carry the harm.
Conditioned on the parser reading the rule, the entire
\imhyp$-$\reactive\ gap concentrates in cells where the
articulated rule is wrong.

\textbf{Qwen-14B: a second failure mode emerges.}
Correct-and-followed cells are nearly at floor (gap $+0.88$), but
\textit{correct\_not\_followed} cells carry a substantial $+3.27$
log-error gap.  This is the \emph{apply-step failure}: the
articulated rule is right, the parser confirms it, the model just
doesn't apply it.  Reactive does not separate articulation from
application and so does not exhibit this mode.

\textbf{Qwen-7B: threshold artefact dominates.}  Even the
\textit{correct\_followed} bin shows a $+106$ gap.  This is not
because correct articulation hurts; it is because at $7$B the
asym-window can include catastrophic arithmetic mistakes on the
$\leq 33\%$ of steps that the cell-level threshold did not
require to be correct.  The 4-way decomposition is therefore not
cleanly diagnostic at this scale; what it shows is that per-step
variability of integer emission overwhelms the rule-articulation
distinction.

\textbf{Coverage caveat.}  The parser is a regex grammar; it
reads roughly $20$--$30\%$ of cells across scales, with the
remainder labelled ``unparseable''.  The unparseable bin still
shows a positive gap to reactive at every scale (gaps $+22.8$,
$+4.5$, $+3.5$, $+2.4$ for $7$/$14$/$32$/$70$B), consistent with
the overall harm pattern.  But these cells cannot be attributed
to either of the two failure modes without additional analysis
(e.g., LLM-as-classifier on the natural-language hypothesis
content).  The headline mechanism claim --- that, conditioned on
parseable rules, the harm concentrates in
\textit{correct\_not\_followed} (at $14$B) and \textit{wrong\_followed}
(at $70$B), neither of which has a reactive analogue --- is
established for the parseable subset, which is sufficient for the
existence claim.

\section{Causal mechanism tests (Benchmark v3)}
\label{sec:bench3}

The 4-way decomposition of the previous section is correlational:
it sorts \imhyp\ cells into bins by what the model happens to
emit, but it does not perturb the prompt to test whether each
bin's harm has the predicted causal source.  Benchmark~v3
introduces four new prompting arms designed to perturb
exactly one factor at a time and so isolate the mechanism.  The
arms reuse v2's deterministic task generator with the same
\verb|--task-seed=2026|, so v3 cells join task-for-task
seed-for-seed with v2.

\paragraph{The four mechanism arms.}
\begin{description}[leftmargin=1.5em,itemsep=2pt,topsep=2pt]
\item[\imhyp{}-SC$_5$:] same prompt as \imhyp, but at each step we
  sample $5$ hypotheses at $T=0.7$ with distinct seeds, parse the
  $5$ integers, and emit the majority vote.  If
  \emph{commitment to a single wrong hypothesis} is the cause, this
  should pull error toward \reactive\ on cells where \imhyp\ is in
  the \textit{wrong\_followed} bin --- because no single wrong
  hypothesis dominates the aggregate.
\item[\imhyp{}-Oracle:] the system prompt is augmented with the
  true generator rule (e.g., $f(k) = (k-1)^2$) and the model is
  instructed to apply it.  This is a sanity check on whether the
  model can apply a correct rule when given one.  If yes,
  \emph{rule-application} works and the harm in \imhyp\ must be in
  \emph{rule-generation}.
\item[\imhyp{}-First:] the prompt is reordered so that the model
  emits the integer FIRST and the hypothesis on the line below.
  If the harm is the textual prefix of the hypothesis conditioning
  the integer-emission step, this should erase the harm.
\item[\imhyp{}-EqBudget:] same prompt as \imhyp\ but with
  \texttt{max\_tokens}\,$=$\,$256$, matching \imcot.  Closes the
  budget-asymmetry caveat (Limitation 5 in Section~\ref{sec:limitations}).
\end{description}

We run these arms on three model scales.  At Qwen-7B and Qwen-14B
we run the full grid ($30$ tasks per band $\times\,3$ seeds, $5$
hypothesis samples for SC$_5$).  At Qwen-32B, where SC$_5$ would
overrun an overnight budget, we run a reduced version with $20$
tasks per band, $2$ seeds, $3$ hypothesis samples, and only the
two most-diagnostic arms (\imhyp{}-SC$_5$ and \imhyp{}-Oracle).
Each run also re-runs the v2 baselines (\reactive, \imhyp,
\imcot) on the same task pool so the analysis is self-contained
and paired.

\subsection{Results}

Table~\ref{tab:bench3} reports the mean log-error gap to
\reactive\ for every (scale, difficulty, arm) cell.  Significance
is from paired $t$-tests on log-error with $n = 90$ paired
observations at $7$B/$14$B and $n = 40$ at $32$B.

\begin{table}[t]
\centering
\caption{Causal mechanism tests on benchmark v3.  Reported is the
mean log-error gap of each arm vs.\ \reactive\ on the same task
pool.  Positive entries hurt; negative entries help.  Bold marks
$p < 0.05$.  Three findings: (1) \imhyp{}-SC$_5$ is statistically
indistinguishable from \imhyp\ at every cell tested
(self-consistency does not help); (2) \imhyp{}-Oracle helps
significantly at $7/9$ cells, with the help \emph{growing with
scale} on hard tasks ($-0.24, -0.25, -0.84$ at $7$/$14$/$32$B);
(3) \imhyp{}-EqBudget is numerically identical to \imhyp\
(token-budget asymmetry was inert).  The arm \imhyp{}-First was
not run at $32$B (overnight-time budget) but at $7$B/$14$B is
indistinguishable from or only mildly milder than \imhyp.}
\label{tab:bench3}
\small
\setlength{\tabcolsep}{4pt}
\begin{tabular}{l l rrrrrr}
\toprule
Model & Difficulty
  & \imhyp & \imcot & SC$_5$ & Oracle & First & EqBudget \\
\midrule
\multirow{3}{*}{Qwen-7B}
  & easy   & $\mathbf{+1.35}^{\hurts}$ & $+0.01$ & $\mathbf{+1.32}^{\hurts}$ & $\mathbf{-0.29}^{\helps}$ & $\mathbf{+1.07}^{\hurts}$ & $\mathbf{+1.35}^{\hurts}$ \\
  & medium & $\mathbf{+1.28}^{\hurts}$ & $+0.06$ & $\mathbf{+1.34}^{\hurts}$ & $\mathbf{-0.30}^{\helps}$ & $\mathbf{+1.48}^{\hurts}$ & $\mathbf{+1.27}^{\hurts}$ \\
  & hard   & $\mathbf{+1.89}^{\hurts}$ & $\mathbf{+0.17}^{\hurts}$ & $\mathbf{+1.58}^{\hurts}$ & $\mathbf{-0.24}^{\helps}$ & $\mathbf{+1.93}^{\hurts}$ & $\mathbf{+1.89}^{\hurts}$ \\
\addlinespace
\multirow{3}{*}{Qwen-14B}
  & easy   & $\mathbf{+1.00}^{\hurts}$ & $\mathbf{+0.71}^{\hurts}$ & $\mathbf{+1.00}^{\hurts}$ & $-0.04$\, & $\mathbf{+0.81}^{\hurts}$ & $\mathbf{+1.00}^{\hurts}$ \\
  & medium & $\mathbf{+1.54}^{\hurts}$ & $\mathbf{+0.89}^{\hurts}$ & $\mathbf{+1.52}^{\hurts}$ & $\mathbf{-0.14}^{\helps}$ & $\mathbf{+0.87}^{\hurts}$ & $\mathbf{+1.54}^{\hurts}$ \\
  & hard   & $\mathbf{+0.99}^{\hurts}$ & $\mathbf{+1.51}^{\hurts}$ & $\mathbf{+0.85}^{\hurts}$ & $\mathbf{-0.25}^{\helps}$ & $\mathbf{+0.95}^{\hurts}$ & $\mathbf{+0.99}^{\hurts}$ \\
\addlinespace
\multirow{3}{*}{Qwen-32B}
  & easy   & $\mathbf{+0.86}^{\hurts}$ & $\mathbf{+0.22}^{\hurts}$ & $\mathbf{+0.78}^{\hurts}$ & $-0.07$\, & --- & --- \\
  & medium & $\mathbf{+1.00}^{\hurts}$ & $\mathbf{+0.50}^{\hurts}$ & $\mathbf{+0.91}^{\hurts}$ & $\mathbf{-0.43}^{\helps}$ & --- & --- \\
  & hard   & $+0.51$\, & $\mathbf{+1.59}^{\hurts}$ & $+0.21$\, & $\mathbf{-0.84}^{\helps}$ & --- & --- \\
\bottomrule
\end{tabular}
\end{table}

\paragraph{Finding 1 (refutation): self-consistency does not help.}
\imhyp{}-SC$_5$ produces gaps essentially identical to \imhyp\ at
every (scale, difficulty) cell tested.  Across the nine cells
tested, the largest gap improvement of SC$_5$ over \imhyp\ is
$+0.31$ log-error (at $7$B-hard); the smallest is $0.0$ (at
$14$B-easy).  None of these improvements approaches \reactive.
The straightforward tokenized-commitment story --- \emph{the
model commits to one wrong hypothesis; sample many and aggregate
to recover} --- is therefore refuted.  The natural reading is
that the population of sampled hypotheses is correlated in its
errors: many sampled hypotheses are wrong in similar ways, and
majority-voting their integers preserves the systematic error
rather than averaging it out.

\paragraph{Finding 2 (positive deployment result):  oracle
helps, and the help grows with scale on hard tasks.}
\imhyp{}-Oracle, in which the true rule is given to the model
directly, helps significantly at $7/9$ cells tested.  The two
non-significant cells (Qwen-14B-easy and Qwen-32B-easy) are
those where \reactive\ is already at the floor (mean error
$0.00$), leaving no room for any arm to improve.  At every cell
where \reactive\ has slack, oracle wins.  The most striking
contrast is on Qwen-32B-hard: \reactive\ scores $20.89$ mean
asymptotic error, oracle scores $3.38$, an $84\%$ reduction
($p = 1.6 \cdot 10^{-3}$).  The benefit of oracle on hard tasks
\emph{grows monotonically with scale}:  $7$B
$-0.24 \to 14$B $-0.25 \to 32$B $-0.84$.  The implication for
practitioners is direct: \emph{when you can give the LLM the
exogenous rule, do so; do not ask it to articulate one.}

\paragraph{Finding 3 (refutation):  textual-prefix conditioning
is not the mechanism.}  \imhyp{}-First reverses the order so the
integer is emitted first; the hypothesis follows on the next line
and cannot condition the integer.  At $7$B all three bands and at
$14$B all three bands, \imhyp{}-First still hurts substantially
($+0.81$ to $+1.93$ log-error gap).  The harm is not primarily
about the autoregressive prefix conditioning the integer.

\paragraph{Finding 4 (cleanup):  the budget-asymmetry caveat is
inert.}  \imhyp{}-EqBudget gives \texttt{max\_tokens}\,$=$\,$256$
to the IM-Hypothesis prompt, matching \imcot.  At every cell at
$7$B and $14$B the result is numerically identical to \imhyp\
within rounding ($\pm 0.001$ log-error).  Limitation~5 of
Section~\ref{sec:limitations} is therefore closed.

\paragraph{Revised mechanism story.}
The four findings together require a sharper account than the
original commitment story.  The harm of the IM-Hypothesis prompt
appears to come from the \emph{rule-generation} step, not from any
of (a) commitment to one sampled rule, (b) the textual ordering
of rule-then-integer, or (c) the token budget.  The model can
apply a stated rule reliably (oracle helps); the model cannot
reliably generate a correct rule, and once it has emitted some
rule, neither resampling nor reordering recovers the implicit
pattern-matcher's accuracy.  This is consistent with the v2 4-way
decomposition (Section~\ref{sec:hyp-correctness}):  cells where
the parser confirms a correct rule are at the reactive floor at
$70$B; cells where the parser confirms a wrong rule carry the
harm.  Benchmark v3 shows that the wrong-rule cells cannot be
fixed by sampling diversity (SC$_5$) or by reordering
(\imhyp{}-First), but they can be eliminated by removing the
generation step entirely (Oracle).

\paragraph{Deployment implication.}
The Oracle result is a genuinely positive finding suitable for
practitioner adoption.  In settings where the rule generating an
exogenous reference is known to the deployer (e.g., a parser is
available, the structure is known a priori, the rule can be
inferred from documentation), the deployer should put the rule
in the system prompt rather than asking the LLM to maintain an
internal model.  The benefit at Qwen-32B-hard is $84\%$ error
reduction; under the rule-generation account, similar benefits
should appear at any scale provided the model can parse the
stated rule.


\section{Cross-paradigm replication: code debugging}
\label{sec:code-debug}

The findings reported in Sections~\ref{sec:results-v2}--\ref{sec:bench3}
establish the Apply-Don't-Articulate phenomenon in deterministic
integer sequence prediction.  A natural reviewer concern is paradigm
specificity: \emph{would the same effects appear in a different
symbolic-reasoning task?}  Code debugging is the natural stress test:
it is symbolic and deployment-realistic; the chain-of-thought
literature reports its largest gains in this regime
\citep{wei2022chain,sprague2024cot}; and grading is fully programmatic
via hidden test suites, eliminating any LLM-as-judge subjectivity.
We apply the v3 condition design verbatim to a code-debugging
benchmark and report a clean cross-paradigm replication of every
v3 mechanism prediction, with effect sizes that are if anything
\emph{larger} than in sequence prediction.

\subsection{Benchmark setup}

We construct $30$ hand-curated buggy Python functions, $10$ per
difficulty band, paralleling the v1 sequence-prediction benchmark
in structure.  Each task supplies (i) a canonical (correct)
implementation, (ii) a buggy implementation differing from the
canonical by a single concrete edit, (iii) a hidden test suite
of four to five pass/fail tests including edge cases, and
(iv) a ground-truth one-line bug description used by the Oracle
arm.

The bands span: \textbf{easy} -- single-token bugs (off-by-one in
\verb|range|, wrong comparison operator, swapped subtraction
operands, missing return, etc.); \textbf{medium} -- control-flow
or data-structure bugs (mutable default argument, append-vs-extend,
accumulator initialised inside the loop, off-by-one in 2D
indexing, etc.); \textbf{hard} -- semantic bugs (wrong algorithm,
wrong sign in distance formula, integer division for float median,
off-by-one in binary search, set-loses-order in dedup, etc.).
A self-validation pass at module import time confirms that for
every template, canonical passes all tests and buggy fails at
least one.  The full template list is in the public repository.

The grader extracts a Python \verb|def| block from the model's
reply (handles markdown fences, prose, multiple defs by selecting
the last syntactically-valid candidate, helper functions),
exec's it in a restricted-builtins namespace (no \verb|import|,
\verb|open|, \verb|exec|), runs the test runner under a SIGALRM
wall-clock timeout, and returns a pass-rate in $[0, 1]$.  The
test runners are in-process callables that are not exposed to
the model; the model never sees the test cases.

\subsection{Conditions}

The seven conditions exactly mirror the v3 design, translated to
the code-debugging idiom.

\begin{itemize}[leftmargin=2em,itemsep=0pt,topsep=2pt]
\item \reactive: ``output the corrected def, no commentary.''
\item \textbf{IM-Diagnosis}: ``state the bug as
  `\texttt{\# BUG: ...}' first, then the def.''  Direct analog of
  IM-Hypothesis.
\item \textbf{IM-CoT}: ``think step by step, then the def.''
\item \textbf{IM-Diag-SC$_5$}: sample $5$ diagnoses at $T = 0.7$,
  group fixes by canonical AST hash, take the modal fix.
\item \textbf{IM-Oracle}: true bug description supplied in the
  system prompt; model just applies it.
\item \textbf{IM-Diag-First}: def first, then
  `\texttt{\# BUG: ...}' comment afterward.
\item \textbf{IM-Diag-EqBudget}: same prompt as IM-Diagnosis with
  \texttt{max\_tokens}\,$=$\,$512$ matching CoT.
\end{itemize}

We test on three scales of the Qwen2.5-Coder family --- the
code-tuned variants, which give the IM prompts the strongest
possible test because the model is trained for the task ---
at $7$B, $14$B, and $32$B.  Per-cell sample size is $n = 30$ at
$7$B and $14$B and $n = 20$ at $32$B (paired by task and seed).

\subsection{Results}

Table~\ref{tab:code-passrate} reports the mean pass-rate per
(scale, difficulty, condition) cell.

\begin{table}[t]
\centering
\caption{Cross-paradigm replication: mean pass-rate per (scale,
difficulty, condition) on the code-debugging benchmark.  Bold
marks cells where the paired $t$-test against \reactive\ on the
arcsin-square-root-transformed pass-rate is significant at
$\alpha = 0.05$ in the hurts direction.  Oracle never
significantly hurts at any cell; mostly ties or trends positive.}
\label{tab:code-passrate}
\small
\setlength{\tabcolsep}{4pt}
\begin{tabular}{l l rrrrrrr}
\toprule
Model & Difficulty
  & \reactive & IM-Diag & CoT & SC$_5$ & Oracle & First & EqBudget \\
\midrule
\multirow{3}{*}{Coder-7B}
  & easy   & $0.975$ & $0.883$              & $0.917$              & $0.925$              & $1.000$ & $0.925$              & $0.883$ \\
  & medium & $0.808$ & $0.825$              & $0.775$              & $0.908$              & $0.917$ & $0.875$              & $0.825$ \\
  & hard   & $0.883$ & $0.850$              & $\mathbf{0.725}^{\hurts}$ & $0.817$              & $0.967$ & $0.850$              & $0.850$ \\
\addlinespace
\multirow{3}{*}{Coder-14B}
  & easy   & $1.000$ & $\mathbf{0.033}^{\hurts}$ & $0.917$              & $\mathbf{0.100}^{\hurts}$ & $1.000$ & $\mathbf{0.867}^{\hurts}$ & $\mathbf{0.033}^{\hurts}$ \\
  & medium & $0.975$ & $\mathbf{0.033}^{\hurts}$ & $0.867$              & $\mathbf{0.083}^{\hurts}$ & $0.975$ & $\mathbf{0.742}^{\hurts}$ & $\mathbf{0.033}^{\hurts}$ \\
  & hard   & $0.967$ & $\mathbf{0.000}^{\hurts}$ & $\mathbf{0.817}^{\hurts}$ & $\mathbf{0.033}^{\hurts}$ & $1.000$ & $\mathbf{0.617}^{\hurts}$ & $\mathbf{0.000}^{\hurts}$ \\
\addlinespace
\multirow{3}{*}{Coder-32B}
  & easy   & $1.000$ & $\mathbf{0.800}^{\hurts}$ & $1.000$              & $1.000$              & $1.000$ & ---                  & --- \\
  & medium & $0.875$ & $\mathbf{0.600}^{\hurts}$ & $0.812$              & $\mathbf{0.650}^{\hurts}$ & $0.975$ & ---                  & --- \\
  & hard   & $0.975$ & $0.800$              & $\mathbf{0.800}^{\hurts}$ & $0.950$              & $1.000$ & ---                  & --- \\
\bottomrule
\end{tabular}
\end{table}

The findings replicate every one of the five v3 mechanism predictions.

\paragraph{Finding 9.1 -- Catastrophic harm at $14$B-Coder.}
Asking Qwen2.5-Coder-14B to prepend a single
`\texttt{\# BUG: ...}' comment line collapses pass-rate from
$\geq 96\%$ to $\leq 3\%$ at every difficulty band.  Paired
$t$-statistics are $25$ to $42$, $p$-values $< 10^{-20}$ at every
cell.  This is the largest single prompt-induced regression
observed in the study, and it occurs on a model trained
specifically for code generation, applied to an explicitly
code-generation task.  The $97$-percentage-point accuracy drop
from a single comment-line addition is, to our knowledge, the
strongest documented case of a prompt-induced capability
regression in any LLM benchmark.

\paragraph{Finding 9.2 -- Inverted-$U$ over scale, replicated.}
The harm of IM-Diagnosis (relative to \reactive) follows the
same scale pattern as the v3 inverted-$U$ for hard-band CoT:
\textbf{barely hurts at $7$B-Coder} (only $1$ of $18$ paired
cells significant), \textbf{catastrophic at $14$B-Coder}
(13/18 cells significant; absolute pass-rate near zero),
\textbf{partial recovery at $32$B-Coder} ($3$ of $9$ cells
significant; absolute pass-rate $0.6$--$0.8$).  The
Coder-7B model largely ignores the diagnosis scaffolding because
the function-fixing pattern is so dominant in its training; the
$14$B model follows the scaffolding literally and produces
wrong fixes under that constraint; the $32$B model partially
recovers, presumably by either sidestepping the scaffolding or
producing more accurate diagnoses.

\paragraph{Finding 9.3 -- Self-consistency does not recover at
the catastrophe scale.}  IM-Diag-SC$_5$ is statistically
indistinguishable from IM-Diagnosis at every $14$B cell
(absolute pass-rate gap $\leq 0.07$ across the three difficulty
bands).  At $32$B, SC$_5$ recovers fully on easy and hard but
not on medium.  The v3 finding that ``self-consistency over
$5$ sampled hypotheses does not fix the simple commitment
story'' replicates verbatim in code: when most sampled
diagnoses are wrong, the modal fix is also wrong, and majority
voting preserves the systematic error.

\paragraph{Finding 9.4 -- Oracle never hurts; trends positive
where slack exists.}  IM-Oracle matches or trends slightly
above \reactive\ at every (scale, difficulty) cell tested.  At
no cell does Oracle significantly hurt at $\alpha = 0.05$.  At
every Coder-14B and Coder-32B cell where \reactive\ has slack
(reactive pass-rate below $1.0$), Oracle's pass-rate is at
least as high.  At Qwen2.5-Coder-32B-medium, Oracle delivers
$+0.10$ pass-rate over reactive ($0.975$ vs.\ $0.875$), and on
$32$B-hard Oracle matches reactive at $1.000$ where reactive
is $0.975$.  The deployment recommendation \emph{provide the
rule rather than ask the model to articulate it} transfers
cleanly from sequence prediction to code debugging.

\paragraph{Finding 9.5 -- Reordering is partially helpful but
not sufficient.}  IM-Diag-First (def first, comment after)
hurts substantially at $14$B-Coder ($\Delta$ pass-rate of
$+0.13$, $+0.23$, $+0.35$ on easy/medium/hard relative to
IM-Diagnosis) but the harm relative to \reactive\ is still
significant at every $14$B band.  At $7$B-Coder the effect is
$n.s.$ across all three bands.  This refines the v3
textual-prefix-conditioning result: at the catastrophe scale,
ordering does provide measurable relief, but the harm of
articulating-then-emitting is largely order-independent.

\paragraph{Finding 9.6 -- Token-budget asymmetry remains inert.}
IM-Diag-EqBudget produces \emph{numerically identical} cells to
IM-Diagnosis at every $14$B cell (pass-rate difference $0.000$
to within rounding).  Limitation~5 of
Section~\ref{sec:limitations} is now closed independently in
both paradigms.

\subsection{Cross-paradigm summary}
\label{sec:cross-paradigm}

Table~\ref{tab:cross-paradigm} consolidates the v3 sequence-prediction
and code-debugging findings.  Every one of the five mechanism
predictions confirmed in v3 transfers to the code-debugging
paradigm.

\begin{table}[t]
\centering
\caption{Cross-paradigm summary of the v3 mechanism predictions.
Each row is a prediction the simple tokenized-commitment account
or one of its variants makes; each column is a paradigm in which
we tested it. \checkmark\ = prediction supported (i.e., empirical
data agrees with the prediction); \texttimes\ = prediction
refuted; ---\ = not tested.}
\label{tab:cross-paradigm}
\small
\begin{tabularx}{\linewidth}{X c c}
\toprule
Mechanism prediction
  & v3 Sequence prediction
  & v3-style Code debugging \\
\midrule
SC over $5$ samples recovers \reactive\ at the catastrophe scale
  & \texttimes\ (refuted)
  & \texttimes\ (refuted) \\
Oracle (rule given) helps significantly
  & \checkmark
  & \checkmark \\
Reordering (def/integer first) erases the harm
  & \texttimes\ (refuted; mild relief only)
  & \texttimes\ (refuted; mild relief only) \\
Token-budget asymmetry explains the harm
  & \texttimes\ (refuted; budget inert)
  & \texttimes\ (refuted; budget inert) \\
Inverted-$U$ harm peaks at intermediate scale
  & \checkmark\ (peak at $14$--$32$B)
  & \checkmark\ (peak at $14$B-Coder) \\
\bottomrule
\end{tabularx}
\end{table}

The Apply-Don't-Articulate phenomenon is now established in two
distinct symbolic paradigms: deterministic integer sequence
prediction (low-vocabulary, single-token output) and code
debugging (high-vocabulary, multi-token output, deployment-
realistic).  All five mechanism predictions transfer cleanly.
The deployment recommendation -- \emph{do not ask the LLM to
articulate its internal model; provide the model with the rule
and ask it to apply it} -- is now backed by data from two
paradigms, seven distinct model scales, and over $10{,}000$
prompted trials.

\section{Discussion}
\label{sec:discussion}

\subsection{What the result is and is not}

The result is that the prompting analogue of the IMP fails:
asking an LLM to externalize an internal model and use it does
not improve closed-loop sequence prediction in the regime tested,
and reliably worsens it.  The result is not that the IMP itself
is wrong; it is a theorem about feedback architectures with
explicit state, not about prompts.  The result also does not
imply that internal models in general are bad; networks
plausibly compute something like an internal model in their
weights and activations, and that implicit model is what
\reactive\ leverages.

\subsection{Comparison with the chain-of-thought literature}

\citet{wei2022chain} demonstrate that CoT helps on complex
multi-step reasoning, particularly for $\geq$ 100B parameter
models.  Our largest tested model is $70$B-q4 -- effectively
smaller after quantization -- so we are formally below the regime
where their strongest gains were observed.  Nevertheless, our
paradigm (deterministic integer-sequence prediction) is symbolic,
the regime where the meta-analysis of \citet{sprague2024cot}
documents the largest CoT gains on average.  Our finding -- that
\imcot\ does not help and that the IMP-flavoured \imhyp\ hurts
substantially across the entire scale range tested -- is therefore
not contradicted by the original CoT literature but does narrow
the regime of CoT applicability further.  In particular, the
inverted-$U$ harm pattern (Section~\ref{sec:cot-inverted-u}) is
not predicted by either of the two main CoT papers and is to our
knowledge new.

\subsection{Implications for treating LLMs as controllers}

The take-away for LLM-as-controller programmes is direct:
the existence of an implicit internal model in an LLM does not
license the prompt-level externalization of that model.  In our
data, the implicit model embedded in network weights and
activations consistently outperforms the model the same network
articulates in tokens.  This has design implications for agentic
systems: prompting an LLM agent to ``maintain a world model'' or
to ``state its plan before acting'' cannot be assumed to
strengthen control performance and may weaken it on
deterministic-tracking tasks.

\subsection{Implications for prompt engineering}

The robustness of the harm under the bootstrap CI (which
resamples tasks) refutes the
``you got unlucky with task choice'' rejoinder.  At three of the
four tested scales, the medium-band IM-Hypothesis harm is
significant under both the parametric paired test and the
non-parametric task-cluster bootstrap.  This is the level of
robustness a careful methodologist requires before reading a
negative result as substantive rather than artifactual.

\subsection{Connection to multi-agent failure modes}

The harm we observe is the single-agent reduction of a class of
failures the MAST taxonomy of \citet{cemri2025why} catalogues at
the system level.  Our benchmark gives a controlled probe for the
single-agent component of those failures, which can then be
composed back to study how an articulated wrong model propagates
when consumed by a downstream agent.  Pre-registered as future
work in Section~\ref{sec:future}.

\section{Limitations}
\label{sec:limitations}

We deliberately use this section to make the scope of our claims
precise.  Each limitation could in principle reverse the
qualitative finding under different choices.

\textbf{Two symbolic paradigms; non-symbolic paradigms untested.}
Section~\ref{sec:results-v2}--\ref{sec:bench3} cover deterministic
integer sequence prediction; Section~\ref{sec:code-debug} covers
Python code debugging.  Both are symbolic-reasoning paradigms.
Continuous-valued tracking, stochastic prediction, dialog state
tracking, multi-step planning, and other non-symbolic deployment
regimes are not tested here.  The cross-paradigm replication
between the two paradigms we do test is consistent with the
generalised claim, but does not establish it for non-symbolic
tasks.

\textbf{Open-weight models only; one model family per paradigm.}
Sequence prediction uses three Qwen2.5 sizes plus one quantized
Llama3.1; code debugging uses three Qwen2.5-Coder sizes.  No
frontier-scale closed model (Claude, GPT, Gemini) is tested in
either paradigm; the cross-paradigm consistency we report is
within the open-weight regime.  Whether the inverted-$U$ peak
shifts at frontier scale, or whether Oracle's monotone help
saturates, are open empirical questions.

\textbf{Three prompt phrasings per paradigm.}  The harm is
specific to the exact prompts in Sections~\ref{sec:prompts}
and~\ref{sec:code-debug}.  Other phrasings of ``maintain an
internal model'' without an explicit
\verb|HYPOTHESIS:|/\verb|# BUG:| line, or with explicit
self-correction instructions, might give different effects.  We
have not run robustness sweeps over the prompt-text dimension.

\textbf{Decoding configuration.}  v1 uses greedy decoding;
v2 and v3 use sampled decoding ($T = 0.7$, $3$--$5$ seeds).  Our
self-consistency arm uses $N = 5$ samples; whether $N \in
\{20, 50\}$ would change the catastrophe-scale result is
untested.  At $N = 5$ the simple commitment story is refuted
in both paradigms; this does not preclude a finer-grained
sampling-diversity account that requires more samples than $5$
to expose.

\textbf{Token-budget asymmetry --- closed.}  \imcot\ originally
used \texttt{max\_tokens}\,$=$\,$256$ versus $96$ for the
hypothesis-prefix arms; we flagged this in earlier drafts as a
potential confound.  The IM-Hyp-EqBudget and IM-Diag-EqBudget
arms (Sections~\ref{sec:bench3} and~\ref{sec:code-debug})
match the budget at $256$ and $512$ tokens respectively, and
produce numerically identical results to their unequal-budget
counterparts in both paradigms.  This caveat is therefore
empirically closed.

\section{Future Work}
\label{sec:future}

The benchmark suite reported here closes several questions that
earlier drafts pre-registered.  We organise this section into
\emph{closed} (questions empirically resolved by the v2/v3/code
benchmarks reported here) and \emph{open} (questions a NeurIPS
reviewer is still entitled to ask).

\subsection*{Closed by this paper}

\textbf{C1.\ Self-consistency over hypotheses.}  Pre-registered
in earlier drafts as F1.  Resolved by Section~\ref{sec:bench3}
(IM-Hyp-SC$_5$ in sequence prediction) and
Section~\ref{sec:code-debug} (IM-Diag-SC$_5$ in code debugging).
Result in both paradigms: $N = 5$ self-consistency does not
recover \reactive\ at the catastrophe scale; refutes the simple
tokenized-commitment account.

\textbf{C2.\ Hypothesis-correctness decomposition.}  Pre-registered
as F2.  Resolved by Section~\ref{sec:hyp-correctness}.

\textbf{C3.\ Reordering / textual-prefix mechanism.}  Pre-registered
as a sub-question of F1.  Resolved by IM-Hyp-First and
IM-Diag-First arms in Sections~\ref{sec:bench3}
and~\ref{sec:code-debug}.  Reordering provides partial relief at
the catastrophe scale but does not erase the harm.

\textbf{C4.\ Token-budget asymmetry.}  Pre-registered as
Limitation~5.  Resolved by IM-Hyp-EqBudget and IM-Diag-EqBudget
arms in both paradigms.  Budget asymmetry is empirically inert.

\textbf{C5.\ Cross-paradigm generality.}  Pre-registered as F6
(continuous-valued targets) but answered at higher value by
Section~\ref{sec:code-debug}'s code-debugging benchmark.  All
five mechanism predictions transfer cleanly between sequence
prediction and code debugging; effect sizes are larger in the
code paradigm.

\subsection*{Open}

\textbf{O1.\ Frontier-model replication.}  Run both benchmarks
on Claude Opus, GPT-5, Gemini, and similar.  If the harm
disappears at frontier scale, the headline tightens to ``below
frontier scale, IMP-as-prompt hurts.''  If it persists in even
one of those three, the claim generalises beyond open-weight
models entirely.  Straightforward but requires API budget that
was out of scope here.  Highest priority.

\textbf{O2.\ Higher-$N$ self-consistency.}  Re-run SC arms at
$N \in \{20, 50, 100\}$ to test whether the simple commitment
story survives at sample sizes where the modal hypothesis
becomes more reliably correct.  At $N = 5$ the story is
refuted; the question is whether a higher-$N$ ``sample over
many wrong rules until the right one wins'' approach can rescue
it.  We expect not, on the basis that the wrong rules are
correlated rather than uniformly random, but this is empirical.

\textbf{O3.\ Non-symbolic paradigms.}  Dialog state tracking,
multi-step planning, real-world agent benchmarks (e.g., SWE-bench).
The cross-paradigm replication between integer sequence
prediction and code debugging is consistent with a general
phenomenon, but both paradigms are symbolic.  A non-symbolic
paradigm where the model has no parseable rule to articulate
might not show the harm at all, or might show it for different
reasons.

\textbf{O4.\ Multi-agent decomposition.}  In a MAS, the IMP
analogue is that one agent maintains an explicit world model
and another plays the controller.  Running the benchmark with a
two-agent (modeller, predictor) split, where the modeller's
output is the predictor's prefix, would test whether the failure
mode persists across agent boundaries (\citealp{cemri2025why}).

\textbf{O5.\ Architectural intervention.}  Fine-tune a small
($7$B) model on data where the integer or fix is emitted
\emph{before} the hypothesis or diagnosis text --- inverting the
standard chain-of-thought ordering.  If the fine-tuned model
loses the rule-generation harm, that is direct architectural
evidence for the cross-paradigm mechanism we propose.

\textbf{O6.\ Mechanistic-interpretability probes.}  Compare the
LLM's internal next-token distribution under \reactive\ to the
distribution under \imhyp\ / IM-Diagnosis conditioned on a
wrong articulated rule.  If the wrong-rule prefix sharply biases
the next-token distribution toward outputs consistent with the
wrong rule, that is direct activation-level evidence for
generation-step commitment.  Combined with the
Section~\ref{sec:bench3} and~\ref{sec:code-debug} behavioural
results, this would identify which finer-grained mechanism
(correlated sampled-hypothesis errors, generation-step
instability) is operative at the activation level.

\section{Conclusion}
\label{sec:conclusion}

We tested the prompting analogue of the Internal Model Principle
in two distinct symbolic-reasoning paradigms: deterministic
integer sequence prediction (over $5{,}760$ prompted trajectories
across four open-weight scales) and Python code debugging (over
$3{,}000$ trials across three Qwen2.5-Coder scales).  Across this
corpus, four findings stand together.

\emph{First}, articulating an internal model never significantly
helps in either paradigm, at any tested scale, under any of the
articulation phrasings tested.

\emph{Second}, the IM-Hypothesis / IM-Diagnosis prompt
significantly hurts in both paradigms.  In sequence prediction
the harm decays monotonically with scale on the easy band
(Spearman $\rho = -1.00$, Pearson $r = -0.998$, $p = 0.002$) but
persists on the medium band at every tested scale (pooled
Stouffer $p < 10^{-6}$).  In code debugging the harm collapses
pass-rate from near-$1.0$ to near-zero at Qwen2.5-Coder-14B at
every difficulty band ($t = 25$--$42$, $p < 10^{-20}$ at every
cell) --- the largest single prompt-induced regression observed
in the entire study.

\emph{Third}, the harm exhibits a scale-dependent inverted-$U$
in both paradigms, peaking at intermediate scale ($14$--$32$B
in sequence prediction; $14$B-Coder in code debugging) and
weakening at the smallest and largest tested scales.  The
mechanism we propose --- \emph{the model articulates
competent-looking but wrong rules in tokens, and once
articulated, the rule constrains downstream emission} --- is
consistent with the inverted-$U$: at small scale articulation is
inert, at intermediate scale articulation is articulate-enough
to lock in but not accurate-enough to be correct, at large scale
articulation is more often correct.

\emph{Fourth}, four causal-mechanism arms confirm the same revised
story in both paradigms.  Self-consistency over $5$ sampled
hypotheses (SC$_5$) does not recover the reactive baseline at
the catastrophe scale, refuting the simplest commitment account.
Reordering (def/integer first) provides partial relief but does
not erase the harm, refuting the textual-prefix account.
Equalising the token budget is empirically inert, closing the
budget caveat.  But \emph{providing} the true rule (the Oracle
arm) helps significantly with growing magnitude on hard tasks,
reaching an $84\%$ error reduction at Qwen-32B-hard in sequence
prediction and $+0.10$ pass-rate at Qwen2.5-Coder-32B-medium in
code debugging.  The revised mechanism story is:  the harm of
the IM-Hypothesis / IM-Diagnosis prompt is in the
\emph{rule-generation} step --- the model can apply a stated
rule reliably but cannot reliably write a correct one.

These results yield a clean deployment finding, not just a
negative one.  In settings where the rule generating an
exogenous reference is known to the deployer (a parser is
available, the structure is documented, the rule can be inferred
from a spec), the deployer should put the rule in the system
prompt rather than ask the LLM to maintain an internal model
of it.  Across two paradigms and seven model scales,
\emph{``maintain an internal model'' prompts hurt;
``apply this stated rule'' prompts help.}

\paragraph{Code and data.}
\url{https://github.com/Minds-R-Lab/phase_margin_llm}.
Discovery benchmark:
\verb|pipeline/experiments/run_imp_benchmark.py|.
Confirmation benchmark:
\verb|pipeline/experiments/run_imp_benchmark2.py|.
Mechanism arms:
\verb|pipeline/experiments/run_imp_benchmark3.py|.
Code-debugging benchmark:
\verb|pipeline/experiments/run_imp_code_benchmark.py| with
\verb|pipeline/src/phase_margin/coding/|.
Per-cell raw outputs and analysis scripts in
\verb|pipeline/results/|.

\bibliographystyle{plainnat}
\begin{thebibliography}{9}

\bibitem[Cemri et al.(2025)]{cemri2025why}
Mert Cemri, Melissa Z. Pan, Shuyi Yang, Lakshya A. Agrawal, Bhavya
Chopra, Rishabh Tiwari, Kurt Keutzer, Aditya Parameswaran, Dan Klein,
Kannan Ramchandran, Matei Zaharia, Joseph E. Gonzalez, and Ion Stoica.
\newblock Why Do Multi-Agent {LLM} Systems Fail?
\newblock arXiv preprint
\href{https://arxiv.org/abs/2503.13657}{arXiv:2503.13657}, 2025.

\bibitem[Francis and Wonham(1976)]{francis1976internal}
B.~A. Francis and W.~M. Wonham.
\newblock The internal model principle of control theory.
\newblock \emph{Automatica}, 12(5):457--465, 1976.
\newblock \href{https://doi.org/10.1016/0005-1098(76)90006-6}%
  {doi:10.1016/0005-1098(76)90006-6}.

\bibitem[Kojima et al.(2022)]{kojima2022zero}
Takeshi Kojima, Shixiang Shane Gu, Machel Reid, Yutaka Matsuo, and
Yusuke Iwasawa.
\newblock Large language models are zero-shot reasoners.
\newblock In \emph{Advances in Neural Information Processing Systems
(NeurIPS)}, 2022.
\newblock arXiv preprint
\href{https://arxiv.org/abs/2205.11916}{arXiv:2205.11916}.

\bibitem[Sprague et al.(2024)]{sprague2024cot}
Zayne Sprague, Fangcong Yin, Juan Diego Rodriguez, Dongwei Jiang, Manya
Wadhwa, Prasann Singhal, Xinyu Zhao, Xi Ye, Kyle Mahowald, and Greg
Durrett.
\newblock To {CoT} or not to {CoT}? Chain-of-thought helps mainly on
math and symbolic reasoning.
\newblock arXiv preprint
\href{https://arxiv.org/abs/2409.12183}{arXiv:2409.12183}, 2024.
\newblock Published at ICLR 2025.

\bibitem[Wang et al.(2023)]{wang2023self}
Xuezhi Wang, Jason Wei, Dale Schuurmans, Quoc Le, Ed Chi, Sharan
Narang, Aakanksha Chowdhery, and Denny Zhou.
\newblock Self-consistency improves chain of thought reasoning in
language models.
\newblock In \emph{International Conference on Learning Representations
(ICLR)}, 2023.
\newblock arXiv preprint
\href{https://arxiv.org/abs/2203.11171}{arXiv:2203.11171}.

\bibitem[Wei et al.(2022)]{wei2022chain}
Jason Wei, Xuezhi Wang, Dale Schuurmans, Maarten Bosma, Brian Ichter,
Fei Xia, Ed Chi, Quoc Le, and Denny Zhou.
\newblock Chain-of-thought prompting elicits reasoning in large
language models.
\newblock In \emph{Advances in Neural Information Processing Systems
(NeurIPS)}, 2022.
\newblock arXiv preprint
\href{https://arxiv.org/abs/2201.11903}{arXiv:2201.11903}.

\end{thebibliography}

\appendix

\section{Discovery-benchmark task list (v1)}
\label{app:tasks}
The 30 hand-curated discovery-benchmark tasks (used in
Section~\ref{sec:method-v1}) are listed in the public repository
under \verb|pipeline/experiments/run_imp_benchmark.py|.

\section{Confirmation-benchmark task families (v2)}
\label{app:families}
The exact parameter ranges, the set of $260$ tasks generated by
\verb|--task-seed=2026|, and the family-by-family per-band counts
are in \verb|pipeline/experiments/run_imp_benchmark2.py| on the
public repository.

\section{Code-debugging task templates}
\label{app:code}
The $30$ buggy-Python-function templates (canonical, buggy,
hidden test suite, ground-truth bug description) are defined in
\verb|pipeline/src/phase_margin/coding/bug_templates.py|.  A
self-validation pass at module-import time confirms that for
every template, canonical passes all tests and buggy fails at
least one.

\section{Selected raw replies (Qwen-7B, Fibonacci, IM-Hyp)}
\label{app:replies}
For illustration of the rule-generation failure mode, the raw
replies on \texttt{h01\_fibonacci} for Qwen-7B under \imhyp\ at
the last six steps of the trajectory ($k = 8, 9, 10, 11, 12, 13$):
true targets are $34, 55, 89, 144, 233, 377$.

\begin{Verbatim}[fontsize=\small,frame=single]
k=8:   HYPOTHESIS: y_k = y_{k-1} + 8     -> 33
k=9:   HYPOTHESIS: y_k = y_{k-1} + 13    -> 47
k=10:  HYPOTHESIS: y_k = y_{k-1} + 13    -> 68
k=11:  HYPOTHESIS: y_k = y_{k-1} + 21    -> 110
k=12:  HYPOTHESIS: y_k = y_{k-1} + 21    -> 165
k=13:  HYPOTHESIS: y_k = y_{k-1} + 34    -> 267
\end{Verbatim}

The model articulates a wrong recurrence with a $d_k$ that
follows the Fibonacci sequence itself but lagged by one step,
and computes the integer prediction deterministically forward
from the lagged $d_k$.  Under \reactive\ on the same task the
model emits the correct integer at every step.  Under
\imhyp{}-Oracle (the rule given in the system prompt), the same
model emits the correct integer at every step.  The failure is
therefore in \emph{generating} the rule, not in \emph{applying}
a known one.

\end{document}
